\section{Methods}\label{sec:methods}

In this section, we present our method for conditioning trajectory diffusion models using compact constraint-response signatures. We first formalize environment-conditioned trajectory generation and describe where conditioning enters the diffusion model. We then introduce diagnostic trajectory snippets, which act as constraint probes for an environment. Finally, we describe an entropy-based greedy procedure for selecting anchor snippets that produce maximally informative environment signatures.

\subsection{Problem Setup}

Let $\mathcal{E}=\{e_1,\ldots,e_N\}$ denote a set of training environments. Each environment specifies a constraint model under which robot trajectories are evaluated. These constraints may include geometric collision constraints, joint-limit constraints, kinematic reachability constraints, torque or dynamic feasibility constraints, contact constraints, or any task-specific validity criterion. We assume access to a constraint oracle
%
\begin{equation}
    \Phi(e,\tau) \in \{0,1\},
\end{equation}
%
where $\Phi(e,\tau)=1$ indicates that trajectory $\tau$ satisfies the constraints of environment $e$, and $\Phi(e,\tau)=0$ indicates that at least one constraint is violated. This abstraction deliberately does not require the constraint oracle to expose the reason for failure. The same representation can therefore be used whether infeasibility is caused by obstacle collision, excessive torque, violation of a velocity bound, or any other constraint encoded by the environment.

We consider trajectories in robot configuration space. A trajectory is represented as
%
\begin{equation}
    \tau = (q_0,q_1,\ldots,q_T), \qquad q_t\in\mathbb{R}^{d},
\end{equation}
%
where $d$ is the number of robot degrees of freedom. The objective is to learn a conditional generative model that samples trajectories from a distribution
%
\begin{equation}
    p(\tau \mid q_{\mathrm{start}}, q_{\mathrm{goal}}, z(e)),
\end{equation}
%
where $q_{\mathrm{start}}$ and $q_{\mathrm{goal}}$ specify the planning query and $z(e)$ is an environment representation used to condition generation. The central question addressed by this work is how to construct $z(e)$. Rather than using a high-dimensional perceptual representation of the scene, we construct a compact behavioral representation that captures how the environment constrains motion.

\subsection{Conditional Trajectory Diffusion}

We model trajectory generation using a denoising diffusion model over trajectories. Given a clean trajectory $\tau_0$, the forward diffusion process gradually corrupts the trajectory with Gaussian noise:
%
\begin{equation}
    q(\tau_t \mid \tau_0)
    =
    \mathcal{N}
    \left(
        \sqrt{\bar{\alpha}_t}\tau_0,
        (1-\bar{\alpha}_t)I
    \right),
\end{equation}
%
where $t\in\{1,\ldots,T_d\}$ is the diffusion timestep and $\bar{\alpha}_t$ is determined by the noise schedule. A neural denoiser $\epsilon_\theta$ is trained to predict the noise added at timestep $t$, conditioned on the planning query and environment representation:
%
\begin{equation}
    \epsilon_\theta
    =
    \epsilon_\theta(
        \tau_t,
        t,
        q_{\mathrm{start}},
        q_{\mathrm{goal}},
        z(e)
    ).
\end{equation}
%
The training objective is the standard noise-prediction loss
%
\begin{equation}
\begin{split}
    \mathcal{L}(\theta)
    =
    \mathbb{E}_{\tau_0,t,\epsilon}
    \big[
    \big\|
        \epsilon -
        \epsilon_\theta(
            \tau_t,
            t,
            q_{\mathrm{start}},
            q_{\mathrm{goal}},
            z(e)
        )
    \big\|_2^2
    \big],
\end{split}
\end{equation}
%
where $\epsilon\sim\mathcal{N}(0,I)$. During inference, trajectories are generated by iterative denoising from random noise while conditioning on the start state, goal state, and environment code.

This formulation places strong importance on the conditioning representation $z(e)$. If $z(e)$ fails to distinguish environments that require different motion behavior, then the denoising process receives insufficient information to generate the appropriate trajectory distribution. Conversely, if $z(e)$ preserves perceptual variation that does not affect feasible motion, then the diffusion model must learn to ignore nuisance variation. We therefore seek a representation that is aligned with behavioral equivalence: environments should be represented similarly when they induce similar feasible trajectory distributions, and differently when they impose different constraints on motion.

\subsection{Diagnostic Trajectory Snippets}

We approximate behavioral equivalence using a finite set of diagnostic trajectory snippets. A snippet is a short contiguous segment extracted from a demonstrated or planned trajectory. Given a trajectory $\tau=(q_0,\ldots,q_T)$, a snippet of horizon $H$ starting at index $a$ is
%
\begin{equation}
    s = (q_a,q_{a+1},\ldots,q_{a+H-1}).
\end{equation}
%
Let $\mathcal{S}=\{s_1,\ldots,s_M\}$ denote the set of candidate snippets extracted from the training trajectories. Each snippet acts as a binary constraint probe. For environment $e_j$, define
%
\begin{equation}
    b_i(e_j) = \Phi(e_j,s_i) \in \{0,1\},
\end{equation}
%
where $b_i(e_j)=1$ indicates that snippet $s_i$ satisfies the constraints in environment $e_j$, and $b_i(e_j)=0$ indicates that it violates at least one constraint. Evaluating all snippets in all training environments produces a binary constraint-response matrix
%
\begin{equation}
    B\in\{0,1\}^{M\times N},
    \qquad
    B_{ij}=b_i(e_j).
\end{equation}
%
Each row of $B$ describes the feasibility pattern of one snippet across the training environments. Each column describes the response of one environment to all candidate snippets. The goal of anchor selection is to choose a subset of snippets whose response patterns are maximally informative for distinguishing environments according to their constraint-induced behavior.

Given an anchor set $A=\{a_1,\ldots,a_K\}\subseteq \{1,\ldots,M\}$, we encode an environment $e_j$ by the binary signature
%
\begin{equation}
    z_A(e_j)
    =
    \left[
        B_{a_1j},
        B_{a_2j},
        \ldots,
        B_{a_Kj}
    \right]
    \in\{0,1\}^{K}.
\end{equation}
%
The signature $z_A(e_j)$ is a behavioral barcode: it records how the environment responds to a fixed set of diagnostic motion probes. Importantly, this code is not intended to reconstruct the environment. It is intended to preserve the distinctions among environments that are relevant to constraint satisfaction and trajectory generation.

\subsection{Entropy-Based Anchor Selection}

A useful anchor is not simply a snippet that is feasible in many environments. A snippet that is feasible in every environment, or infeasible in every environment, provides little information for conditioning because it does not separate environments. Instead, an anchor is informative when its constraint-response pattern refines the current grouping of environments into more distinguishable signature classes.

Let $A_m$ denote the set of anchors selected after $m$ greedy iterations, and let
%
\begin{equation}
    C_m(e_j)=z_{A_m}(e_j)
\end{equation}
%
be the current signature assigned to environment $e_j$. The current signatures induce a partition of the training environments into equivalence classes
%
\begin{equation}
    \mathcal{G}_m = \{G_c : c \in \mathrm{range}(C_m)\},
\end{equation}
%
where
%
\begin{equation}
    G_c = \{e_j \in \mathcal{E} : C_m(e_j)=c\}.
\end{equation}
%
Within each group $G_c$, all environments are indistinguishable under the anchors selected so far. A candidate snippet is informative if it splits these existing groups into smaller subgroups.

We formalize this using conditional entropy. For a candidate snippet $s_i\notin A_m$, define the empirical probability that the snippet is feasible within group $G_c$:
%
\begin{equation}
    p_{i,c}
    =
    \frac{1}{|G_c|}
    \sum_{e_j\in G_c} B_{ij}.
\end{equation}
%
The uncertainty of the candidate response within group $G_c$ is measured by binary entropy:
%
\begin{equation}
    h(p_{i,c})
    =
    -p_{i,c}\log_2 p_{i,c}
    -(1-p_{i,c})\log_2(1-p_{i,c}),
\end{equation}
%
with $h(0)=h(1)=0$. The entropy gain of candidate snippet $s_i$ conditioned on the current anchor set is
%
\begin{equation}
    \Delta H(s_i\mid A_m)
    =
    \sum_{G_c\in\mathcal{G}_m}
    \frac{|G_c|}{N}
    h(p_{i,c}).
\end{equation}
%
This score is high when a snippet produces uncertain, balanced feasibility outcomes inside large groups of environments that are currently indistinguishable. It is zero if the snippet has the same outcome for every environment within each current group, since such a snippet would not refine the existing signatures. At each iteration, we select the candidate snippet with maximum conditional entropy:
%
\begin{equation}
    a_{m+1}
    =
    \arg\max_{i\notin A_m}
    \Delta H(s_i\mid A_m).
\end{equation}
%
The selected snippet is appended to the anchor set,
%
\begin{equation}
    A_{m+1}=A_m\cup\{a_{m+1}\},
\end{equation}
%
and every environment signature is updated by appending the corresponding feasibility bit. The procedure terminates when $K$ anchors have been selected, when the maximum entropy gain falls below a threshold, or when the current signatures are sufficiently discriminative for the downstream task.

This objective can be interpreted as greedily maximizing the information provided by the next constraint probe about the environment identity, conditioned on the probes already selected. Equivalently, the method chooses the snippet whose binary response is most uncertain within the current signature classes. As a result, anchors are selected to expose constraint-relevant variation that has not yet been captured by earlier probes.

\paragraph{Anchor selection procedure.}
Given training environments $\mathcal{E}=\{e_1,\ldots,e_N\}$, candidate snippets $\mathcal{S}=\{s_1,\ldots,s_M\}$, a constraint oracle $\Phi$, and a desired maximum number of anchors $K$, we construct the anchor set as follows:
Rather than presenting anchor identification as a planner-specific subroutine, we view it as a representation construction procedure. The procedure greedily refines a probe-induced partition of the training environments.

\begin{enumerate}
    \item \textbf{Construct constraint responses.}
    Evaluate every candidate snippet in every training environment to form the binary constraint-response matrix
    %
    \begin{equation}
        B_{ij} = \Phi(e_j,s_i),
        \qquad
        B\in\{0,1\}^{M\times N}.
    \end{equation}
    %
    Each row of $B$ corresponds to the feasibility pattern of a snippet across environments.

    \item \textbf{Initialize signatures.}
    Set the anchor set to $A=\emptyset$ and assign every environment the empty signature,
    %
    \begin{equation}
        z_A(e_j) = [\ ],
        \qquad \forall e_j\in\mathcal{E}.
    \end{equation}

    \item \textbf{Form current equivalence classes.}
    At greedy iteration $m$, group environments by their current signatures. This induces a partition
    %
    \begin{equation}
        \mathcal{G}_m = \{G_c : c \in \mathrm{range}(z_A)\},
    \end{equation}
    %
    where
    %
    \begin{equation}
        G_c = \{e_j\in\mathcal{E}: z_A(e_j)=c\}.
    \end{equation}
    %
    Environments in the same group are indistinguishable under the anchors selected so far.

    \item \textbf{Score candidate anchors.}
    For each unselected snippet $s_i\notin A$, compute the empirical probability that the snippet is feasible within each current group:
    %
    \begin{equation}
        p_{i,c}
        =
        \frac{1}{|G_c|}
        \sum_{e_j\in G_c} B_{ij}.
    \end{equation}
    %
    The conditional entropy gain of snippet $s_i$ is then
    %
    \begin{equation}
        \Delta H(s_i\mid A)
        =
        \sum_{G_c\in\mathcal{G}_m}
        \frac{|G_c|}{N}
        h(p_{i,c}),
    \end{equation}
    %
    where
    %
    \begin{equation}
        h(p)=-p\log_2 p-(1-p)\log_2(1-p),
    \end{equation}
    %
    with $h(0)=h(1)=0$.

    \item \textbf{Select the most informative snippet.}
    Choose the candidate snippet with maximum conditional entropy gain:
    %
    \begin{equation}
        a^\star
        =
        \arg\max_{s_i\notin A}
        \Delta H(s_i\mid A).
    \end{equation}
    %
    This selects the snippet whose constraint-response pattern most strongly refines the current environment partition.

    \item \textbf{Update anchor set and signatures.}
    Add the selected snippet to the anchor set,
    %
    \begin{equation}
        A \leftarrow A\cup\{a^\star\},
    \end{equation}
    %
    and append its response bit to every environment signature:
    %
    \begin{equation}
        z_A(e_j)
        \leftarrow
        [z_A(e_j), B_{a^\star j}],
        \qquad \forall e_j\in\mathcal{E}.
    \end{equation}

    \item \textbf{Terminate.}
    Repeat Steps 3--6 until $K$ anchors have been selected or the maximum conditional entropy gain falls below a threshold $\eta$.
\end{enumerate}

% \begin{algorithm}[t]
% \caption{Entropy-Based Diagnostic Anchor Selection}
% \label{alg:anchor_selection_abstract}
% \begin{algorithmic}[1]
% \Require Training environments $\mathcal{E}=\{e_1,\ldots,e_N\}$, candidate snippets $\mathcal{S}=\{s_1,\ldots,s_M\}$, constraint oracle $\Phi$, maximum anchors $K$, entropy threshold $\eta$
% \Ensure Anchor set $A$ and environment signatures $\{z_A(e_j)\}_{j=1}^{N}$

% \State Construct the constraint-response matrix
% \[
%     B_{ij} = \Phi(e_j,s_i), \qquad
%     B\in\{0,1\}^{M\times N}.
% \]

% \State Initialize the anchor set $A\gets\emptyset$ and assign each environment the empty signature:
% \[
%     z_A(e_j) \gets [\ ], \qquad \forall e_j\in\mathcal{E}.
% \]

% \For{$m=0$ to $K-1$}

%     \State Let $\mathcal{G}_m$ be the partition of environments induced by the current signatures:
%     \[
%         e_i,e_j \in G_c
%         \quad \Longleftrightarrow \quad
%         z_A(e_i)=z_A(e_j)=c.
%     \]

%     \State For each unselected candidate snippet $s_i\notin A$, compute its conditional entropy gain:
%     \[
%         \Delta H(s_i\mid A)
%         =
%         \sum_{G_c\in\mathcal{G}_m}
%         \frac{|G_c|}{N}
%         h\!\left(
%             \frac{1}{|G_c|}
%             \sum_{e_j\in G_c} B_{ij}
%         \right),
%     \]
%     where
%     \[
%         h(p)=-p\log_2 p-(1-p)\log_2(1-p).
%     \]

%     \State Select the most informative snippet:
%     \[
%         a^\star
%         =
%         \arg\max_{s_i\notin A}
%         \Delta H(s_i\mid A).
%     \]

%     \If{$\Delta H(a^\star\mid A)<\eta$}
%         \State \textbf{break}
%     \EndIf

%     \State Add the selected snippet to the anchor set:
%     \[
%         A\gets A\cup\{a^\star\}.
%     \]

%     \State Update each environment signature by appending the selected response bit:
%     \[
%         z_A(e_j)
%         \gets
%         \left[
%             z_A(e_j), B_{a^\star j}
%         \right],
%         \qquad \forall e_j\in\mathcal{E}.
%     \]

% \EndFor

% \State \Return $A,\{z_A(e_j)\}_{j=1}^{N}$

% \end{algorithmic}
% \end{algorithm}

% \begin{algorithm}[t]
% \caption{Entropy-Based Diagnostic Anchor Selection}
% \label{alg:anchor_selection}
% \begin{algorithmic}[1]
% \Require Training environments $\mathcal{E}=\{e_1,\ldots,e_N\}$, candidate snippets $\mathcal{S}=\{s_1,\ldots,s_M\}$, constraint oracle $\Phi$, maximum anchors $K$, entropy threshold $\eta$
% \Ensure Anchor set $A$, environment signatures $\{z_A(e_j)\}_{j=1}^{N}$

% \State Initialize $B\in\{0,1\}^{M\times N}$
% \For{$i=1$ to $M$}
%     \For{$j=1$ to $N$}
%         \State $B_{ij} \gets \Phi(e_j,s_i)$
%     \EndFor
% \EndFor

% \State $A \gets \emptyset$
% \For{$j=1$ to $N$}
%     \State $z_A(e_j) \gets [\ ]$
% \EndFor

% \For{$m=0$ to $K-1$}
%     \State Construct groups $\mathcal{G}_m$ induced by current signatures $\{z_A(e_j)\}_{j=1}^{N}$
%     \State $\Delta H^\star \gets -\infty$
%     \State $a^\star \gets \mathrm{None}$

%     \For{$i=1$ to $M$}
%         \If{$i\in A$}
%             \State \textbf{continue}
%         \EndIf

%         \State $\Delta H_i \gets 0$

%         \ForAll{$G_c\in\mathcal{G}_m$}
%             \State $p_{i,c} \gets \frac{1}{|G_c|}\sum_{e_j\in G_c} B_{ij}$
%             \If{$p_{i,c}=0$ or $p_{i,c}=1$}
%                 \State $h_{i,c} \gets 0$
%             \Else
%                 \State $h_{i,c} \gets -p_{i,c}\log_2 p_{i,c}-(1-p_{i,c})\log_2(1-p_{i,c})$
%             \EndIf
%             \State $\Delta H_i \gets \Delta H_i + \frac{|G_c|}{N}h_{i,c}$
%         \EndFor

%         \If{$\Delta H_i > \Delta H^\star$}
%             \State $\Delta H^\star \gets \Delta H_i$
%             \State $a^\star \gets i$
%         \EndIf
%     \EndFor

%     \If{$a^\star=\mathrm{None}$ or $\Delta H^\star < \eta$}
%         \State \textbf{break}
%     \EndIf

%     \State $A \gets A\cup\{a^\star\}$
%     \For{$j=1$ to $N$}
%         \State $z_A(e_j) \gets [z_A(e_j), B_{a^\star j}]$
%     \EndFor
% \EndFor

% \State \Return $A,\{z_A(e_j)\}_{j=1}^{N}$
% \end{algorithmic}
% \end{algorithm}

\subsection{Probe-Induced Behavioral Equivalence}

The selected anchor set defines a finite approximation to behavioral equivalence. Two environments $e_i$ and $e_j$ are equivalent under anchor set $A$ if they produce the same constraint-response signature:
%
\begin{equation}
    e_i \sim_A e_j
    \quad\Longleftrightarrow\quad
    z_A(e_i)=z_A(e_j).
\end{equation}
%
This equivalence relation does not require the environments to be perceptually or geometrically similar. Instead, it groups environments according to their response to diagnostic motion probes. Environments with the same signature may differ in appearance or raw geometry, but they are indistinguishable with respect to the selected constraint queries. Conversely, environments that are perceptually similar may be assigned different signatures if they impose different feasibility outcomes on the anchors.

This representation is therefore aligned with the needs of trajectory generation. The diffusion model does not require a complete description of the environment; it requires a conditioning signal that distinguishes environments when their feasible trajectory distributions differ. The anchor signature serves as a compact index into this space of constraint-induced motion behavior.

\subsection{Encoding New Environments}

After the anchor set $A$ has been selected, encoding a new environment $e_{\mathrm{new}}$ requires no retraining of the anchor selection procedure. We evaluate the same anchor snippets under the constraint oracle of the new environment:
%
\begin{equation}
    z_A(e_{\mathrm{new}})
    =
    \left[
        \Phi(e_{\mathrm{new}},s_{a_1}),
        \ldots,
        \Phi(e_{\mathrm{new}},s_{a_K})
    \right].
\end{equation}
%
The resulting binary vector has the same dimension and coordinate semantics as the training signatures. It can therefore be used directly as the environment-conditioning input for the trajectory diffusion model. If the selected anchors capture recurring modes of constraint variation in the training distribution, then new environments with similar constraint-response structure will map to nearby signatures, even when their raw perceptual descriptions differ.

The same procedure applies regardless of the underlying constraint type. For collision-only planning, $\Phi$ may check whether each configuration in the snippet is collision-free. For dynamic planning, $\Phi$ may additionally check velocity, acceleration, torque, or contact constraints. For task-constrained motion, $\Phi$ may encode semantic or operational requirements. In all cases, the anchor selection algorithm operates only on the binary constraint-response matrix and is therefore agnostic to the implementation of the constraint oracle.

\subsection{Algorithm Summary}

The full anchor construction procedure is summarized in Algorithm~\ref{alg:anchor_selection}. The method first extracts candidate snippets from feasible training trajectories and evaluates every snippet in every training environment to form the constraint-response matrix $B$. It then initializes all environments with the empty signature and greedily adds snippets that maximize the conditional entropy of the candidate response within current signature groups. The resulting signatures provide compact environment codes for diffusion conditioning. Unlike sensory representations, these codes are explicitly constructed to preserve constraint-induced behavioral distinctions rather than perceptual detail.
